\chapter{CONCLUSION} \label{c7}
\justifying

\section{Conclusion}

This project has successfully developed a Real-Time Sign Language Interpreter system that translates American Sign Language (ASL) gestures into text and speech in real-time. The system achieves the objectives set forth at the beginning of this project.

The key achievements of this project include:

\begin{enumerate}
    \item \textbf{Successful Implementation}: Built a complete sign language recognition system using MediaPipe for hand landmark detection and Random Forest for gesture classification.
    \item \textbf{Real-Time Performance}: The system processes video at 16 FPS with 53ms total latency, well within the 100ms target.
    \item \textbf{Hardware Accessibility}: Works on standard laptop CPU without GPU requirements, making it accessible to a wide range of users.
    \item \textbf{Complete Features}: Includes text output, text-to-speech, stability checking, and intuitive web interface.
\end{enumerate}

The system successfully addresses the communication barrier between deaf and hearing communities by providing an automated, real-time translation of sign language gestures. The implementation demonstrates that computer vision and machine learning can be effectively applied to solve real-world accessibility problems.

The technical approach used - combining MediaPipe's efficient hand tracking with a Random Forest classifier - provides an optimal balance between accuracy, speed, and computational requirements. This aligns with the findings from the literature survey, which showed similar approaches achieving 93-98\% accuracy.

All functional and non-functional requirements specified in the requirements analysis have been met. The system passes all unit tests, integration tests, and performance benchmarks. User acceptance testing shows positive feedback with a 4.2/5 overall satisfaction rating.

This research contributes to the field of assistive technology by demonstrating an accessible, real-time sign language recognition system that can be deployed on consumer hardware without specialized equipment.

\vspace{1em}
\section{Future Work}

While the current system meets all initial requirements, there are several directions for future enhancement.

\subsection{Word-Level Recognition}
The current system recognizes individual letters. A future enhancement would involve implementing LSTM/GRU networks to enable word-level recognition, which would allow more natural translation and faster input. The primary challenge lies in the requirement for video sequence processing rather than static frame analysis.

\subsection{Multi-Hand Support}
At present, the system handles only a single hand. Future work could extend this to detect both hands simultaneously, enabling the recognition of two-handed signs. However, this would introduce increased complexity in the processing pipeline.

\subsection{Cloud Deployment}
The current implementation is limited to local deployment. Migrating to a cloud-based API would significantly improve scalability and accessibility for a broader user base. Key challenges to address in this transition include managing latency and ensuring user privacy.